\documentclass[11pt,answers]{exam}
\usepackage[legalpaper,bindingoffset=0in,left=.6in,right=0.8in,top=0.4in,bottom=0.7in,headsep=.5\baselineskip]{geometry}
\usepackage[utf8]{inputenc}
\usepackage{graphicx}
\usepackage{listings}
\usepackage{enumitem}
\usepackage{amsmath}
\usepackage{multirow}
\usepackage{multicol}
\usepackage{caption}
\usepackage{tabularx}
\usepackage{url}
\usepackage{amssymb}
\usepackage{array}
\usepackage{tikz}
\usetikzlibrary{shapes.callouts, positioning, arrows.meta, shapes.geometric, shadows, calc, patterns, 3d, backgrounds, shadings}
\usepackage{adjustbox}
\usepackage{ragged2e}
\usepackage{pgfplots}
\pgfplotsset{compat=1.18}

\def\changemargin#1#2{\list{}{\rightmargin#2\leftmargin#1}
\item[]}
\let\endchangemargin=\endlist

\usepackage{xcolor}
\definecolor{mGreen}{rgb}{0,0.6,0}
\definecolor{mGray}{rgb}{0.5,0.5,0.5}
\definecolor{mPurple}{rgb}{0.58,0,0.82}
\definecolor{backgroundColour}{rgb}{0.95,0.95,0.92}
\lstset{
  language=C,
  basicstyle=\ttfamily\small,
  aboveskip={1.0\baselineskip},
  belowskip={1.0\baselineskip},
  columns=fixed,
  extendedchars=true,
  breaklines=true,
  tabsize=4,
  prebreak=\raisebox{0ex}[0ex][0ex]{\ensuremath{\hookleftarrow}},
  frame=lines,
  showtabs=false,
  showspaces=false,
  showstringspaces=false,
  keywordstyle=\color[rgb]{0.627,0.126,0.941},
  commentstyle=\color[rgb]{0.133,0.545,0.133},
  stringstyle=\color[rgb]{01,0,0},
  numbers=left,
  numberstyle=\small,
  stepnumber=1,
  numbersep=10pt,
  captionpos=t,
escapeinside={\%*}{*)}
}

\pointsdroppedatright

\title{
\vspace{-1.5cm}
\begin{changemargin}{0.68cm}{0.68cm}
  \normalsize Class Test 4 Makeup\hfill {July 2025}\\
\end{changemargin}
\vspace{.4cm}
\large BANGLADESH UNIVERSITY OF ENGINEERING AND TECHNOLOGY\\
\vspace{.1cm}
Department of Computer Science and Engineering\\
\vspace{.2cm}
Course Code: CSE 219 \hspace{.1cm} Course Title: Signals and Linear Systems\\
\vspace{.2cm}
\textbf{Time: 25 minutes} \hspace{1.7cm} \textbf{Marks: 20}
\begin{changemargin}{0.68cm}{0.68cm}
  \normalsize Student Name: \underline{\hspace{6.5cm}} \hfill Student ID: \underline{\hspace{2.5cm}}\\
\end{changemargin}
}
\date{}

\begin{document}

\maketitle
\thispagestyle{empty}
\vspace{-2.4cm}

\begin{center}
Answer \textbf{all of Questions 1--3}. Answer concisely.
\end{center}

\begin{questions}

% Question 1
\question Consider the sequence
\[
  x[n] = \{2,0,1,-1\}, \qquad 0 \le n < 4
\]
You want to compute its 4-point DFT using the \textbf{radix-2 DIF FFT}. \hfill [8]

\begin{parts}
  \part Perform the first butterfly stage and write the intermediate values.
  \part Complete the second stage and write the outputs in the order produced directly by the DIF algorithm.
  \part Rearrange the result from part (b) to obtain the DFT in natural order, i.e. $X[0], X[1], X[2], X[3]$.
  \part Briefly state where bit-reversal appears in radix-2 DIT FFT and where it appears in radix-2 DIF FFT.
\end{parts}
\vspace{0.2cm}

% Question 2
\question In an audio-processing application, a discrete-time signal $x[n]$ of length $L$ is passed through a filter $h[n]$ of length $M$ to produce
\[
  y[n] = (x*h)[n].
\]
For large $L$ and $M$, you want to compute this efficiently. \hfill [8]

\begin{parts}
  \part What is the length of the output sequence $y[n]$? What is the time complexity of direct time-domain computation?
  \part If you compute DFTs of insufficient length, multiply them, and then take an IDFT, what kind of convolution do you obtain instead of linear convolution?
  \part Propose an FFT-based algorithm to compute the correct filtered output.
  \part If the transform length is $N$, what condition must $N$ satisfy? What is the resulting asymptotic complexity?
\end{parts}
\vspace{0.2cm}

% Question 3
\question[4] Thank you for attending the classes and participating in this CT. \hfill [4]

\end{questions}

\end{document}
